{\Large\bfseries Predicting the Next Category and Region of a Visit}\par
\smallskip
{\large Speaker cards v3: 13 main slides, about 16 minutes of talk, about 13 minutes for questions}\par
\smallskip
Slot: 30 minutes \emph{including} questions. First talk of the symposium, right after the 8:15
opening. Audience: wireless networking and mobility management, not machine learning.
Deck: \texttt{slides/main.pdf}. Footer numbers 1--12 are the main slides; backups B1--B12 follow.

\medskip
\begin{itemize}
  \item \textbf{Time marks}: end of S2b $\approx$ 3:20 $\cdot$ end of S4 $\approx$ 6:20 $\cdot$ end of S7
        $\approx$ 11:15 $\cdot$ end of S9 $\approx$ 13:25 $\cdot$ end $\approx$ 16:30.
  \item \textbf{If short of time}: trim S3 and S5, never S2b, S7 or S10.
  \item \textbf{Three sentences never cut}: the 2-point margin fixed in advance (S4); ``deficits, not
        ties'' (S7); ``one pass costs more compute; the gain is operational'' (S10).
  \item \textbf{Jump to a backup}: B$k$ is PDF page $13+k$ (B1 = 14 \dots{} B12 = 25). Type the
        number, then Enter.
\end{itemize}

% ---------------------------------------------------------------
\card{S1 $\cdot$ Title}{20 s}{title, three authors, NESPeD-LAB / UFV.}
\say
\begin{itemize}
  \item ``Good morning. I am Vitor Hugo, from the Federal University of Viçosa, in Brazil. This is
        joint work with Germano dos Santos and Fabrício Silva.''
  \item ``We ask whether one model can predict both what kind of place a person visits next, and in
        which neighborhood.''
\end{itemize}

% ---------------------------------------------------------------
\card{S2 $\cdot$ One person's day, and two questions}{90 s}{schematic map, two boxes, a definition
line.}
\textbf{\color{alert}Rehearse aloud. It wins or loses the room at 8:30.}\par
\say
\begin{itemize}
  \item ``On a location-based social network such as Gowalla, people share the places they visit.
        Each shared visit is a check-in: a place, which this literature calls a point of interest, and
        a time.''
  \item ``Here is one person's day: nine check-ins. Coffee, gym, office, lunch. The tenth is the
        question.''
  \item ``A large body of work predicts the exact next place; Texas alone has about a hundred and
        sixty thousand of them. We predict two properties of the next visit instead.''
  \item ``What kind of place: one of seven types, such as food, shopping, nightlife. And which
        neighborhood: a census tract in the U.S., a \emph{mahalle} in Istanbul; from 520 in Istanbul to
        8,501 in California.''
  \item ``Both are things a service can act on. The kind of place says what the user will want. The
        neighborhood says where to prepare: cache content where the user is heading, provision
        capacity before demand arrives. Over many users, that is where demand is heading.''
  \item ``Handover prediction already enables anticipatory adaptation at the radio cell; Vielhaus and
        colleagues showed it here at MobiWac in 2022. Our unit is a different one: the neighborhood.
        A census tract is not a cell, and we claim nothing about cells or handover.''
\end{itemize}
\ifasked{``Is that a real user?'' No. It is a schematic of the input format, not a case from the
data. ``Why the tenth?'' It is simply the next visit: the model reads the last nine visits and predicts
the one that follows. Nine is the window length, fixed during development and not tuned on results.}
\never{``nobody needs the next place''; ``easier''; ``most work''; a tract is ``coarser'' than a cell
(the paper never compares sizes).}

% ---------------------------------------------------------------
\card{S2b $\cdot$ Two models, or one?}{90 s}{two options, the gap line, the question in a box, the
short answer.}
\say
\begin{itemize}
  \item ``The reference point is one dedicated model per question: two things to train, version and
        serve.''
  \item ``Multi-task learning puts both in one model: one artifact to train, version and deploy, and
        the two tasks might help each other. But sharing has a cost: the shared parameters can settle on
        a compromise that is good for neither task, helping one while hurting the other.''
  \item ``The field increasingly models several granularities at once, but in those systems, such as
        MCMG and HMT-GRN, category and region are helpers for predicting the exact next place. To our
        knowledge, predicting a fine-grained region as an answer in its own right, next to category and
        with no next-place target, is underexplored.''
  \item ``So this is our question: can one model predict both, and what does it cost?''
  \item ``The short answer, so you can decide whether to keep listening: yes, at a measured price. On
        the neighborhood question, one model outperforms a dedicated model at the two largest datasets
        and is at most 0.87 points behind at the other four. On the kind-of-place question, the two are
        within half a point, and how we represent a visit mattered more than one model versus two. Here
        is how.''
\end{itemize}
\ifasked{``Do MCMG and HMT-GRN use multi-task learning?'' Yes, both. HMT-GRN learns the next place
and the next regions jointly, with the next place as its main task; MCMG adds next-region and
next-category heads to a next-place model as auxiliary channels. They are multi-task, but category and
region are there to help the next-place answer; here they are the answers.}
\never{``known to fail''; ``first''; ``few multi-task models''; ``one cost'' or ``cheaper''.}

% ---------------------------------------------------------------
\card{S3 $\cdot$ Check2HGI: a fourth level below the place}{100 s}{one line on place embeddings, the
Check2HGI diagram, two lines.}
\say
\begin{itemize}
  \item ``A model cannot read a place name; it needs a vector per input. Place embeddings learn one
        vector per place, so that similar or nearby places sit close together. Deep Graph Infomax learns
        them on a network of places. Hierarchical Graph Infomax adds the geography: place, then region,
        then city. Both give one vector per place. So two visits to the same coffee shop, Tuesday
        morning and Saturday night, look identical to the model.''
  \item ``We add a fourth level below the place: the check-in itself. Each visit gets its own
        64-dimensional vector, built from its place's category, the hour and day of week, and four
        elapsed-time values, such as the time since the previous visit.''
  \item \emph{[point left to right]} ``Reading the figure left to right: each check-in becomes a
        vector; those vectors are pooled into their place; and the place, region and city levels sit on
        top. This per-visit vector is what our model reads. The magenta copy is a shuffled graph, used
        only as a training signal.''
  \item ``Two things to keep. The graph is trained without the task labels: it never sees the next
        category or the next region. And it is built from the past alone: edges between consecutive
        visits run forward in time only.''
\end{itemize}
\ifasked{``Is it pre-trained?'' The per-visit vectors come from an unsupervised graph trained on the
same data; it is not an external pre-trained encoder. ``What are the four elapsed-time values?'' The
time since the user's previous visit and since their first visit (both log-scaled), the gap within the
same day, and a first-visit flag. Each is measured up to the visit itself, never after it.}

% ---------------------------------------------------------------
\card{S4 $\cdot$ How it was measured}{80 s}{five rows: Regions, Split, Scored, Repeated, Decided.}
\say
\begin{itemize}
  \item ``Before the results, how we measured. Regions are census tracts, or \emph{mahalle} in
        Istanbul: 520 in Istanbul to 8,501 in California.''
  \item ``We split by user: the prediction models never train on a test user's visits. The
        representation is the exception, and I come back to it in the limits.''
  \item ``Category is scored with macro-F1, which counts the seven categories equally; always
        predicting the most common one scores between 5.7 and 7.3. Region is scored with Acc@10: is the
        true region among the model's top ten? A random top ten is right at most about two percent of
        the time.''
  \item ``Each configuration is trained twenty times: four seeds, five folds each.''
  \item \textbf{Never cut:} ``And for region we wrote down what counts as not worse, two Acc@10 points,
        before we saw any result. Two points is below the level at which a service that acts on which
        neighborhood will be busy would behave differently.''
\end{itemize}

% ---------------------------------------------------------------
\card{S5 $\cdot$ Result 1: the representation}{100 s}{six-row table, check-in level vs place level
(HGI).}
\say
\begin{itemize}
  \item ``First result, and the most controlled one: only the input changes. Same model, same folds,
        same windows.''
  \item ``With a vector per visit instead of per place, next-category prediction improves at all six
        datasets and in every fold, by 0.23 to 2.81 points.''
  \item ``The gain is larger at the three smaller datasets and under one point at the three largest.
        At Florida the fold-to-fold spread is larger than the gap, and the paired test does not separate
        the two.''
  \item \textbf{Say this before anyone asks:} ``And to be clear: a control that simply adds the raw
        per-visit features to the place vector gains as much or more at Alabama and Florida, and less at
        Arizona; we ran it at three datasets. What the graph adds beyond those features, this control
        does not separate.'' (Backup B9.)
\end{itemize}
\nums{AL $+1.62$, AZ $+2.58$, Istanbul $+2.81$ $\cdot$ Florida $+0.23$, fold sd $0.42$, $p = 0.067$.}
\ifasked{``These numbers differ slightly from slide 7.'' This table is one seed; slide 7 averages four.
CA and TX place cells: their row matching could not be checked from what is on disk; do not defend them
as verified pairs. Both gains are under one point.}

% ---------------------------------------------------------------
\card{S6 $\cdot$ One model for two tasks}{95 s}{the model figure; four short lines on the right.}
\say
\begin{itemize}
  \item ``Now the one model. In multi-task learning, one model does several tasks at once by sharing
        most of its parts. The risk: whatever the two tasks share can settle on a compromise worse for
        each than a dedicated model would be.''
  \item ``Here is how this model is built. Each task has its own input and its own small encoder, with
        no shared weights. The category task reads the window of per-visit vectors. The region task
        reads the same visits, each represented by the vector of its region.''
  \item ``Both go into the shared trunk, the only place where the two tasks meet. In the trunk,
        attention lets each task read the other's features, while each keeps its own weights. So the
        tasks share by exchanging information, not by owning layers in common.''
  \item ``The region task also keeps a private spatial path that the category task never touches: it
        keeps what is purely spatial out of the exchange.''
  \item ``Whether one model holds up against two is the next slide.''
\end{itemize}
\ifasked{``What happens in the bidirectional cross-attention?'' Each side reads the other: the category
side picks what is useful from the region side, and the region side does the same back. See the
study notes at the end for the plain version.}
\never{count blocks or widths; ``the private path was added because of negative transfer'' (untested).}

% ---------------------------------------------------------------
\card{S7 $\cdot$ One model, two answers: the price}{100 s}{Dedicated vs Joint; green = outperforms;
bold = the larger.}
\say
\begin{itemize}
  \item ``This is the price of one model instead of two.''
  \item ``On region, the joint model outperforms the dedicated model at the two largest datasets, Texas
        and California, by about one point, in twenty of twenty folds. At the other four it is behind by
        at most 0.87 points, inside the two-point margin we fixed in advance.''
  \item ``Those four are deficits, not ties: each interval lies below zero, Istanbul only just. The
        intervals are in a backup slide.''
  \item ``On category, the joint model is ahead at Florida. The other five differences are small, in
        both directions, and unresolved; every interval is within half a point.''
  \item \textbf{The thesis sentence:} ``The trade is a measured one and not a free substitution.''
  \item ``Two cautions, with the numbers. The joint model is also the larger model, and this design
        does not separate size from sharing. And California has more regions than Texas and a smaller
        gain, so that pattern is an observation, not a law.''
\end{itemize}
\nums{region: Texas $+1.21$, California $+1.06$ (20 of 20 folds); largest deficit Alabama $-0.87$
$\cdot$ category: Florida $+0.19$ (19 of 20 folds).}
\ifasked{``Your $\Delta$ is not the difference of the columns.'' $\Delta$ is the paired mean over the
four seeds, so it can differ from the rounded columns in the second decimal.}
\never{``ties'', ``matches'', ``no difference'', ``the same'' for the four region cells.}

% ---------------------------------------------------------------
\card{S8 $\cdot$ And against published systems}{70 s}{four externals vs Dedicated vs Joint; bold best,
underline second; a four-line legend.}
\say
\begin{itemize}
  \item ``Four published systems, all run on our data. POI-RGNN predicts the next category; we
        re-implemented it. HMT-GRN is a multi-task model for next place and next regions: we kept its
        multi-task structure, removed the next-place parts, and ran it on our windows and folds, read at
        its best epoch like our models. STAN is a next-location attention model, re-implemented with its
        own representations, and adapted to rank regions. ReHDM is a reference under its own published
        split.''
  \item ``Both our models, the dedicated and the joint, are above every system we compared against;
        the joint model by at least 4.11 points on category and 3.55 on region. That margin includes the
        representation's advantage, so it is not evidence for sharing.''
\end{itemize}
\ifasked{``Do they run on your folds?'' Only HMT-GRN runs on our windows and folds; STAN and ReHDM use
their own splits. The 3.55 is against STAN at Florida on its own split; on our folds the closest is
HMT-GRN at Florida, 3.78.}
\never{``both our models by at least four points''. The margins are the joint model's only.}

% ---------------------------------------------------------------
\card{S9 $\cdot$ What a service would receive}{60 s}{three labelled lines and a caveat box.}
\say
\begin{itemize}
  \item ``What would a service receive? First, time to act: half of the next visits come within under
        half an hour in Florida, and within five and a half hours in Istanbul; five to twenty-seven
        percent come more than three days later.''
  \item ``Second, a shortlist. Ask the model for its ten most likely neighborhoods. In California, with
        eight and a half thousand neighborhoods, the true one is in that list sixty-four percent of the
        time; in Texas, sixty-six. That is the Acc@10 you saw on slide seven, read as a hit rate.''
  \item ``For scale, a transition table with no learning already gets fifty-one to seventy-two; the
        model is four to ten points above it.''
  \item ``We built and evaluated no such service. This is motivation, not a measured result.''
\end{itemize}
\never{any distance in kilometers; ``24 minutes'' (the median is about 26: say ``under half an
hour'').}

% ---------------------------------------------------------------
\card{S10 $\cdot$ What it costs, and four limits}{95 s}{parameters and the compute sentence (left),
four limits (right).}
\say
\begin{itemize}
  \item ``What does one model cost? The joint model is larger than the two dedicated models combined:
        4.2 million parameters against 1.9 million at Alabama, 5.2 against 2.8 at California, and one
        forward pass costs more compute than running the two small models. The gain is operational, not
        compute: one artifact to train, version and deploy, and two answers.''
  \item Limit 1: ``The representation is trained once over all users. Rebuilding it per fold from
        training users only moved the dedicated models by at most 0.17 macro-F1 and 0.13 Acc@10, at one
        seed and at the datasets checked.''
  \item Limit 2: ``Epoch selection reads the fold the score is read on, so absolute scores are
        optimistic. The joint-versus-dedicated comparison is affected far less.''
  \item Limit 3: ``The vector of a visit is built from the visits before it and nothing after. That
        costs context a two-way model would have. It is also what keeps any visit vector from carrying
        its own target, so I would not trade it.''
  \item Limit 4: ``The joint model is larger, so the two region gains may come from size rather than
        from sharing. This design does not separate them, and I do not credit sharing for them. What
        does not depend on that is the result itself: one model, both answers, at these levels.''
\end{itemize}
\ifasked{latency or energy: ``About five million parameters and a window of nine visits per query. We
did not measure latency or energy.''}

% ---------------------------------------------------------------
\card{S11 $\cdot$ Conclusion}{70 s}{Built / Measured / Not shown / Next, and the tagline.}
\say
\begin{itemize}
  \item ``To conclude. We built two things: a check-in-level representation that extends hierarchical
        graph infomax from the place to the individual visit, and one model that answers both
        questions.''
  \item ``What we measured: on region the joint model outperforms the dedicated model at Texas and
        California and stays inside the two-point margin elsewhere, as deficits, not ties; on category
        the two are within half a point. The trade is measured, not free.''
  \item ``What this does not show is that sharing is the cause: on category the representation
        matters more than one model versus two, and the region gains may come from size.''
  \item ``The results are promising on the point a service cares about, one model carrying both
        answers at a measured price, and they leave room to work: a search for the joint model at Texas
        and California, separating size from sharing, and measuring how far a miss lands, inside a
        service that uses the shortlist.''
  \item Close: ``What a user will do next, and where, from one model instead of two. Thank you.''
\end{itemize}
\never{``a new approach''; ``unified model''; ``in one forward pass instead of two'' (the joint pass
costs more compute).}

% ---------------------------------------------------------------
\card{S12 $\cdot$ Thank you}{20 s}{Questions? Code, data, funders.}
\say
\begin{itemize}
  \item ``The code, the model, the representation, the baselines and the statistical tests are public
        at this link, and both datasets are public. This work was supported by FAPEMIG, CNPq and CAPES.
        I am happy to take questions.''
  \item \textbf{Leave this slide up during questions.} Jump to a backup only when needed, then come back.
\end{itemize}

% ===============================================================
\newpage
\section{Questions: the prepared answers}

\qa{``Your numbers do not match the paper.'' \emph{(most likely: the PDF is in the proceedings)}}{``The paper mixed sources from different versions; the corrected values come from the delivered
build.'' For the HMT-GRN column: ``the paper read HMT-GRN at a fixed final epoch; here it is read at its
best validation epoch, the same rule as our models.'' That correction \emph{favors} the competitor.
Example if asked: ``the Istanbul place-level cell was computed on an older build; matched on the same
rows the gain is 2.81.'' The message survives: a gain at all six datasets and every fold. Do not
apologize or dramatize.}{S8}
\qa{``Isn't it just a bigger model?'' \emph{(this room will ask)}}{First, the paper-level answer: ``Possibly. The design does not separate size from sharing, so we
declare it rather than claim transfer. The claim is about the model: one artifact gives both answers,
inside the margin, and above the dedicated model at the two largest datasets.'' \textbf{Only if pressed}
(``did you test it?''): ``One single-seed control at California, not in the paper: a dedicated region
model at the joint model's parameter count scores 0.41 above it, so at California the gain measures
capacity. Texas has no matched control.'' Never deny it was run; never put it on a slide.}{S10}
\qa{``How do you know a visit vector does not see its target?''}{By the mechanism: each visit node draws only on the visits before it, in training and at readout, and
the graph is trained without the task labels.}{B4}
\qa{``What do MCMG and HMT-GRN read: check-ins or places?''}{A sequence of check-ins, but each check-in is represented by a vector learned for its place ID, one vector per place. MCMG adds learned vectors for the region (k-means clusters) and the category, plus time, distance and position; our HMT-GRN run is an LSTM over learned place-ID vectors. Ours gives each visit its own vector, learned from the check-in graph.}{S2b}
\qa{``Do the external systems run on your folds?''}{``Only HMT-GRN runs on our windows and folds. STAN and ReHDM use their own splits.''}{B7}
\qa{``Where does the 2-point margin come from? Is that standard?''}{Not a convention of this field. Non-inferiority margins come from clinical trials, where the margin is
fixed in advance from what matters in practice (TOST, Lakens 2017). Ours was written before any result:
two points is below the level at which a service acting on which neighborhood will be busy would behave
differently. The seed-to-seed standard deviation of the paired difference is 0.02 to 0.16, far below
the margin.}{B2}
\qa{``Why not one point?''}{``At one point, Istanbul, Arizona and Florida would still pass; Alabama would not.''}{B10}
\qa{``Would a Markov chain do this?''}{A first-order Markov chain over regions gets 51.23 to 72.47 Acc@10; the joint model exceeds it by 4.1 to
10.0 at every dataset. Say the narrowest first: Florida, $+4.07$. Talk only about \emph{our} margin
over the floor; at its best epoch HMT-GRN is above the floor too.}{B3}
\qa{``Isn't this just the timestamp added to the place vector?'' \emph{(the most dangerous one)}}{Largely, yes, on category: a control that adds the raw features to the place vector gains as much or
more at Alabama and Florida, less at Arizona ($+1.73$ / $+1.70$ / $+1.02$). What the graph adds beyond
that, this control does not separate.}{B9}
\qa{``Why not give the vector its neighbors, or later visits?''}{Nearby places are already linked at the place level. What is excluded on purpose is the user's
\emph{later} visits: letting them in is exactly what would let a vector carry its target.}{B4}
\qa{``Does it work for new places or new users?''}{Not without retraining: the representation is trained on each dataset's graph (transductive). An
inductive variant that embeds unseen places and users is future work.}{none}
\qa{``Then what is multi-task learning for?''}{On category the representation moves the score more than the choice between one model and two. Whether
the exchange between tasks adds anything beyond it is not separated here. What one model gives for
certain is operational.}{B6}
\qa{``Was the tuning fair?''}{The dedicated category model received the wider search; the joint model at Texas and California uses a
configuration transferred from the other datasets. Where the searches differ, it favors the dedicated
model.}{none}
\qa{``What does one model instead of two cost?''}{4.2 M vs 1.9 M parameters (AL), 5.2 M vs 2.8 M (CA), against the two dedicated models combined; one
pass costs more compute.}{B5}
\qa{``Does the gain grow with the number of regions?''}{No. California has more regions than Texas and a smaller gain. An observation, not a law.}{B8}
\qa{``And the confidence intervals?''}{The twelve cells with 90\% intervals.}{B10}
\qa{``What are the datasets?''}{Five Gowalla states and Istanbul (Massive-STEPS): counts, regions, windows, lengths, density and the
majority share are in the backup table.}{B12}
\qa{``Why census tracts and not a grid?''}{A grid cell is the usual target; we use official neighborhood-scale units, 520 to 8,501 per dataset.}{B1}
\qa{``How far off are the misses, in kilometers?'' \emph{(rehearse: it comes as a friendly question)}}{``Not measured. It needs per-visit predictions the evaluation path does not keep. It is the first
thing we would measure for a service.'' \textbf{Never give a distance.}}{none}
\qa{``Latency? Energy? Can it run on the edge?''}{``About five million parameters and a window of nine visits per query. We did not measure latency or
energy.''}{none}
\qa{``You split by user, not by time. Drift? Retraining?''}{``The split is by user and not by time, so the evaluation does not measure drift or the effect of
sporadic events.'' The upside: the prediction models never saw the test users.}{none}
\qa{``Gowalla is from 2009--2011.''}{Istanbul comes from Massive-STEPS, a recent collection, and Gowalla is the standard benchmark of this
line. What matters is the pattern that repeats across continents and periods.}{none}
\qa{``Privacy: the representation was trained on everyone.''}{Server side, on the operator's own logs; no privacy mechanism is claimed. Limit 1 measures the effect:
rebuilding it with training users only moved the dedicated models by at most 0.17 macro-F1 and 0.13
Acc@10, at one seed and the datasets checked.}{S10}
\qa{``Only four runs?''}{Each seed is a full five-fold cross-validation: twenty trained models per configuration. The standard
deviation of the paired difference is 0.02 to 0.16 across the four seeds.}{S4}
\qa{``Who would consume this?''}{Name the consumer: load and demand anticipation at neighborhood level, caching, capacity planning, by an
operator or a city service.}{S9}


\section{Never say}
\begin{itemize}
  \item ``beats'' or ``wins''. The verb is \textbf{outperforms}.
  \item ``first'', ``a new approach'', ``few multi-task models''. The paper says the pair as end targets
        is \emph{underexplored}, and its novelty is the combination.
  \item ``known to fail''. Sharing \emph{can} cost; this paper measures it for this pair.
  \item ``ties'', ``matches'', ``no difference'' for the four region cells. They are deficits inside a
        registered margin.
  \item ``in one forward pass instead of two'' or ``cheaper''. The joint pass costs more compute.
  \item Any distance in kilometers. Handover or cell-level gains as a claim.
  \item ``Both our models by at least four points''. The margins on S8 belong to the joint model.
  \item ``The Markov floor beats the external systems''.
  \item ``The gain comes from sharing''.
  \item Any number read from the printed paper where the deck shows a different one. Use the slide.
\end{itemize}

\section{On the day}
\begin{itemize}
  \item 8:30 in Paris is 4:30 in Brasília (France leaves summer time on 25 Oct).
  \item Arrive before the 8:15 opening and test the projector with the PDF. The deck is 16:9.
  \item Carry the PDF on a USB stick and in the cloud.
  \item The program publishes no room or chair. Introduce yourself to the chair before the session.
  \item Leave slide 12 (code, data, funders) on screen during questions.
\end{itemize}

% ===============================================================
\newpage
\section{Study notes: your questions, in plain words}
\textit{For understanding, not for reading aloud. Each answer is checked against the paper or the
dissertation, chapter 5.}

\study{Is the check-in level above or below the place?}
Both readings are true. In the hierarchy, the check-in is the finest level: city, region, place, and
then the check-in at the bottom. In the data flow, everything starts from it: each check-in becomes a
vector, the vectors of one place are pooled into that place, and the place, region and city levels sit
on top. The vector the model reads is the per-visit one.

\study{Why embeddings at all, and why DGI and HGI?}
A model cannot read ``Starbucks on 5th Street''; it needs numbers. An embedding is a learned list of
numbers per item, placed so that similar or nearby items are close. DGI learns one per place on a
network of places. HGI adds the geography: place, region, city. Both give one vector per place, so a
Tuesday-morning visit and a Saturday-night visit to the same café look the same. Check2HGI gives each
visit its own vector.

\study{What are the four elapsed-time values?}
Time since the user's previous visit, time since the user's first visit (both on a log scale), the gap
within the same day, and a flag for ``first visit''. Each is measured up to the visit itself, never
after it.

\study{What happens in the bidirectional cross-attention?}
Picture each task as a person holding nine notes, one per visit in the window.
\begin{enumerate}
  \item The category person looks at the region person's nine notes and asks: which of these help me?
        They take a weighted mix of the useful ones and add it to their own notes.
  \item At the same time, the region person does the same with the category notes. Reading in both
        directions is what ``bidirectional'' means.
  \item Each person then tidies their own notes with their own small layer (the feed-forward layer).
  \item This happens twice: two blocks.
\end{enumerate}
Each direction has its own weights. What crosses between the tasks is information, not layers.

\study{So what do the two tasks actually share?}
The information they read from each other in the trunk. Because of that reading, the category error
also trains the region side, and the other way round. They do not own hidden layers in common.

\study{Why the ``tenth'' visit? Why nine?}
There is nothing special about ``tenth'': the model reads the last nine visits and predicts the next
one. Nine is the window length. It was fixed during development and not tuned on the results.

\study{Do MCMG and HMT-GRN use multi-task learning? What do they read?}
Yes, both are multi-task. HMT-GRN learns the next place and the next regions together, with the next
place as its main task. MCMG adds next-region and next-category outputs to a next-place model, as
``auxiliary channels''. Both read a sequence of check-ins, but represent each one by a vector learned
for its place ID (MCMG also adds region, category and time vectors). The difference here: category and
region are the answers, not helpers, and each visit has its own vector.

\study{Why a 2-point margin? Is that common? Why not 1?}
It is not a convention of this field. Non-inferiority margins come from clinical trials: you fix, in
advance, the smallest difference that would matter in practice, and then test that you are not worse
than that (TOST, Lakens 2017). Ours was written before any result: two Acc@10 points is below the level
at which a service that acts on which neighborhood will be busy would behave differently. The noise is
much smaller: the paired difference varies 0.02 to 0.16 across seeds. With a 1-point margin, Istanbul,
Arizona and Florida would still pass, Alabama would not.

\study{Is limit 3 really a limit? Isn't it what prevents leakage?}
It is both, and the paper says both in the same sentence. A visit's vector uses nothing after the
visit, so it has less context than a model that reads the whole sequence both ways. That same rule is
what stops a vector from carrying its own target. Note: nearby places are already linked in the graph;
what is left out on purpose is the user's later visits.

\study{Doesn't limit 4 kill the shared model?}
No. The paper never claims that sharing works; it claims that one model gives both answers at a
measured price, and that holds whatever causes the region gains. And the limit is real: in one
single-seed control at California, a dedicated region model with 97\% of the joint model's parameters
scored 0.41 Acc@10 above it (5 of 5 folds, $p = 0.010$). So at California the region gain measures
capacity. Texas has no matched control. Say this only if pressed, and never put it on a slide.

\study{What do the numbers on slide 9 mean?}
\emph{Time to act}: half of the next visits come within 0.4 hours (about 26 minutes) in Florida and
5.5 hours in Istanbul; 5 to 27\% come more than 3 days later. \emph{Shortlist}: if a service prepares
the model's top ten neighborhoods, the right one is in that list 64.54\% of the time in California
(out of 8,501) and 66.15\% in Texas. That is simply Acc@10 read as a hit rate. \emph{Scale}: a
transition table with no learning already reaches 51 to 72; the model is 4 to 10 points above it.

\study{How were the external systems run?}
POI-RGNN: re-implemented, next category. HMT-GRN: its multi-task structure kept, the next-place parts
removed, run on our windows and folds, one seed, read at its best epoch; not the complete published
system. STAN: re-implemented with its own representations and splits, adapted to rank regions. ReHDM:
its own published split. Only HMT-GRN uses our folds.

\study{Why do slides 5 and 7 differ in the second decimal?}
Slide 5 is one seed (five folds). Slide 7 averages four seeds. And on slide 7 the $\Delta$ is the
paired mean over the seeds, so it can differ from the difference of the rounded columns.

\study{Why not say ``first'' or ``a new approach''?}
Other works already give one vector per visit (CTLE) or predict category and region inside a
next-place model. What is new is the combination: per-visit context inside a hierarchical graph
representation, and region as an answer of its own next to category. The paper's word is
``underexplored''.

